\setcounter{section}{3}

\Needspace{5\baselineskip}
\hypertarget{attention-residuals}{%
\section{Attention Residuals}\label{attention-residuals}}

\begin{quote}
本文是论文阅读笔记，内容代表对应论文方法或作者理解，不应直接视为领域共识或工程最佳实践。
\end{quote}

\Needspace{5\baselineskip}
\hypertarget{ux4e00ux95eeux9898ux80ccux666f}{%
\subsection{一、问题背景}\label{ux4e00ux95eeux9898ux80ccux666f}}

主流大模型内部的残差连接在每一层处理完信息之后，把处理结果和原来的输入加在一起，传给下一层。每一层做完累加之后，信息就被压进了一个混合状态里。越往后，这个混合状态就越臃肿。随着网络越来越深，混合状态的数值越来越大，发生PreNorm稀释：越深的层，贡献越被稀释。实验证明，一个很深的模型里，删掉相当多的层，效果几乎不变。我们实际上浪费了大量算力。

\Needspace{5\baselineskip}
\hypertarget{ux4e8cux6838ux5fc3ux601dux60f3-1}{%
\subsection{二、核心思想}\label{ux4e8cux6838ux5fc3ux601dux60f3-1}}

目前的主流架构中，每一层的输入=网络输入+前面所有层的输出（等权重）。Kimi团队提出，可以让每一层=网络输入+前面所有层的输出（按需加权求和），权重由这一层自己学。

具体而言，就是用注意力机制，让模型自主决定关注哪些层。每一层有一个自己专属的小向量（只有一个，极其轻量），用它去跟前面所有层的输出算相似度，相似度高的层权重就大，相似度低的权重就小，最后加权求和。值得注意的是，注意力层和MLP层可以关注不同的历史层。

这样，早期层的信息就不会被淹没，可以按需取用。整个改动增加的参数量，相比几十亿参数的大模型来说可以忽略不计，但性能提升到原来的1.25倍，相对于同等性能下节省20％算力。

\Needspace{5\baselineskip}
\hypertarget{ux4e09block-attention-residuals}{%
\subsection{三、Block Attention Residuals}\label{ux4e09block-attention-residuals}}

这种方法的工程问题在于，每一层要关注前面所有层的输出，就必须把前面所有层的输出都存着。现代大模型一般有上百层，训练时还会把模型切成很多份分布到不同机器上（流水线并行），每台机器都要维护所有层的输出，内存压力和通信压力都会暴增。

Kimi的解法是，不用存每一层，改为存每一组。把所有层分成若干个块（block），每个块内部还是用传统方式累加，但把整个块压缩成一个向量存起来。跨块之间，再用注意力机制选择。这样存储量从正比于层数，变成正比于块数。实验发现，分成大约8块就能恢复Full Attention Residuals绝大部分的效果，而内存和通信开销只是原来的一小部分。

\FloatBarrier
\Needspace{5\baselineskip}
\hypertarget{ux53c2ux8003ux6587ux732e-104}{%
\subsection{参考文献}\label{ux53c2ux8003ux6587ux732e-104}}

\vspace{0.25em}
\begin{itemize}
\tightlist
\item
  Kimi Team. (2026). \href{https://arxiv.org/abs/2603.15031}{Attention Residuals}. arXiv:2603.15031.
\end{itemize}

\clearpage
